% ══════════════════════════════════════════════════════════════════════════════
%  EXPERIMENT 9 — THE GREAT GREEN MIX UP                         (Lesson 18, Term 2 Week 7)
% ══════════════════════════════════════════════════════════════════════════════
% NOTE FOR KODIE: written without the faculty GGMU sheet, as a general
% "plan your own separation" practical. The mixture below (split peas, sand,
% iron filings, salt, green food colouring) is an assumption: swap it for the
% faculty mixture and the tasks still work.
\experimentdiv{The Great Green Mix Up}{Planning a sequence of separations \textbullet{} Writing a method in the third person \textbullet{} Judging how well it worked}{Lesson 18 \textbullet\ Term 2, Week 7 \textbullet\ Lab}\label{sec:9}
\lessoncodes{VC2S8U06, U07 \textbullet{} VC2S8I01--I03, I06, I08}

\vocabhead
\begin{ctab}
\begin{tabular}{|SW{4.2cm}|W{11.6cm}|}
  \hline\hd & \thead{Word} & \thead{Definition (copy from the slide)} \\ \hline
  & \scaf{\textbf{Mixture}} & \scaf{Two or more substances mixed together but not chemically joined.} \\ \hline
  & \textbf{Property} & \\ \hline
  & \textbf{Sieving} & \\ \hline
  & \textbf{Magnetic separation} & \\ \hline
  & \textbf{Flow chart} & \\ \hline
\end{tabular}
\end{ctab}

\begin{lipanel}
\inquiryli{Plan, carry out and evaluate a sequence of separations; write the method in the third person.}
\understandli{Choose a separation technique for each substance based on its properties.}
\sctitle
\begin{itemize}[leftmargin=1.4em, itemsep=0pt, topsep=1pt]
  \item I can match each substance to a property that lets me separate it.
  \item I can put the separations in an order that works, and show it as a flow chart.
  \item I can say how pure each separated substance was.
\end{itemize}
\end{lipanel}

\begin{skillsfocus}
\focus{QP2}{Write the aim of the challenge.}
\focus{PC1}{Plan the order of the separations.}
\focus{PC3}{Plan the safety of every step.}
\focus{PC7}{Choose equipment for each separation.}
\focus{CO1}{Write the method in the third person.}
\focus{EV3}{Judge how well each separation worked.}
\end{skillsfocus}

% ─── TECHNIQUE ────────────────────────────────────────────────────────────────
\exppart{Technique spotlight}{Blue}

\begin{technique}{Choosing a separation by property}{PC7}
Every separation uses a \textbf{property} that one substance has and the others do not.
\begin{ctab}
\renewcommand{\arraystretch}{1.3}
\begin{tabular}{|L{4.4cm}|L{4.6cm}|L{5.4cm}|}
  \hline\thead{Property} & \thead{Technique} & \thead{Example} \\ \hline
  Size of the pieces & Sieving & Gravel from sand \\ \hline
  Magnetic & Magnetic separation & Steel cans from aluminium cans \\ \hline
  Insoluble in water & Filtration, decanting & Sand from salt water \\ \hline
  Soluble in water & Dissolving, then evaporation & Salt from sand \\ \hline
  Different boiling points & Distillation & Pure water from salt water \\ \hline
  Different solubility in a solvent & Chromatography & The dyes in green ink \\ \hline
\end{tabular}
\end{ctab}
\end{technique}

\begin{workedex}{Planning the order}
\textbf{Mixture:} sand, salt and paperclips.
\begin{enumerate}[leftmargin=1.6em, itemsep=0pt, topsep=2pt]
  \item Magnetic separation removes the paperclips (steel is magnetic).
  \item Water is added and stirred: the salt dissolves, the sand does not.
  \item Filtration: the sand is the residue, the salt water is the filtrate.
  \item Evaporation of the filtrate leaves the salt.
\end{enumerate}
The magnet goes first: once the water is added, wet sand would stick to the paperclips.
\end{workedex}

\begin{wedo}{Planning the order}
A mixture contains rice, flour and salt. Plan the order of separations and give the property each step uses.
\answerbox{5cm}
\end{wedo}

% ─── PLAN ────────────────────────────────────────────────────────────────────
\exppart{Plan the separation}{Teal}

Your teacher gives your group a beaker of the \textbf{Great Green Mix Up}: green split peas, sand, iron filings and salt, all dyed with green food colouring. Your challenge is to get back as much of each as you can, as pure as you can.

\task Write the aim of the challenge. \skr{QP2}
\answerbox{2.5cm}

\needspace{12\baselineskip}
\task Complete the table. \skr{PC7}
\begin{ctab}
\begin{tabular}{|SW{3.4cm}|W{4.8cm}|W{3.6cm}|W{3.2cm}|}
  \hline\hd & \thead{Substance} & \thead{Property I will use} & \thead{Technique} & \thead{Equipment} \\ \hline
  & \scaf{Split peas} & \scaf{Much larger pieces} & \scaf{Sieving} & \scaf{Kitchen sieve} \\ \hline
  & Iron filings & & & \\ \hline
  & Sand & & & \\ \hline
  & Salt & & & \\ \hline
  & Green colouring & & & \\ \hline
\end{tabular}
\end{ctab}

\needspace{20\baselineskip}
\task Draw a flow chart of your separations, from the whole mixture at the top to the separate substances at the bottom. Write the technique on each arrow. \skr{PC1}
\drawbox{9cm}{}

\needspace{12\baselineskip}
\task Complete the risk assessment. \skr{PC3}
\begin{ctab}
\begin{tabular}{|SW{4.6cm}|W{5.0cm}|W{5.4cm}|}
  \hline\hd & \thead{Hazard} & \thead{Risk} & \thead{Control} \\ \hline
  & \scaf{Iron filings} & \scaf{Could get in eyes} & \scaf{Safety glasses; keep the magnet in a plastic bag} \\ \hline
  & & & \\ \hline
  & & & \\ \hline
\end{tabular}
\end{ctab}

{\small Teacher's initials before you start: \rule{2.5cm}{0.4pt}}

% ─── RESULTS ─────────────────────────────────────────────────────────────────
\exppart{Results}{Amber}

\task Write the method your group actually followed, as numbered steps in the third person. \skr{CO1}
\begin{ctab}
\begin{tabular}{|SW{1.4cm}|W{13.6cm}|}
  \hline\hd & \thead{Step} & \thead{What was done} \\ \hline
  & \scaf{1} & \scaf{The mixture was poured through a sieve and the peas were collected on paper towel.} \\ \hline
  & 2 & \\ \hline
  & 3 & \\ \hline
  & 4 & \\ \hline
  & 5 & \\ \hline
  & 6 & \\ \hline
  & 7 & \\ \hline
\end{tabular}
\end{ctab}

\needspace{12\baselineskip}
\subsubsection*{Table 9.1: How pure was each substance?}
\begin{ctab}
\begin{tabular}{|SW{3.0cm}|C{3.2cm}|W{8.6cm}|}
  \hline\hd & \thead{Substance} & \thead{Purity (1 to 5)} & \thead{What was still mixed in, or lost} \\ \hline
  & \scaf{Split peas} & \scaf{4} & \scaf{Still stained green; a little sand stuck to them} \\ \hline
  & Split peas & & \\ \hline
  & Iron filings & & \\ \hline
  & Sand & & \\ \hline
  & Salt & & \\ \hline
\end{tabular}
\end{ctab}

% ─── ANALYSIS ────────────────────────────────────────────────────────────────
\exppart{Analysis and evaluation}{Coral}

\task Which separation worked best? Explain why, using the property it relied on.
\answerbox{4cm}

\needspace{10\baselineskip}
\task Which substance was least pure? Name what was still mixed in, and how it got there. \skr{EV3}
\answerbox{4cm}

\needspace{10\baselineskip}
\task Did the order you planned work? If you changed it during the lesson, explain why.
\answerbox{4cm}

\needspace{10\baselineskip}
\task The green colouring dissolved in the water, so it ended up with the salt. Suggest how the salt could have been made white again, or how the colouring could be separated from the water.
\answerbox{4cm}

\needspace{10\baselineskip}
\task Suggest two improvements to your method and explain how each would make a substance purer. \skr{EV3}
\answerbox{5cm}

\begin{yr8box}
\begin{keepext}
\exttask Recycling plants separate mixed household waste. Choose three materials that a plant must separate, and for each give the property and technique it could use.
\answerbox{5cm}
\end{keepext}
\end{yr8box}

% ─── SKILL CHECK ─────────────────────────────────────────────────────────────
\exppart{Skill check}{NavyBlue}

\begin{skillcheck}
\scrow{QP2}{I can write the aim of a challenge.}
\scrow{PC1}{I planned the separations in an order that works.}
\scrow{PC3}{I planned how to make each step safe.}
\scrow{PC7}{I chose equipment that suits each separation.}
\scrow{CO1}{I wrote the method in the third person.}
\scrow{EV3}{I can judge how pure each substance was and suggest improvements.}
\end{skillcheck}

\task Look back over all nine experiments. Which skill have you improved most this unit? Give an example.
\answerbox{5cm}

\end{document}
